% =============================================================================
%  E05 — Informe de Dashboard v1  |  Sprint 3  |  SAMO
%  Proyecto: Análisis y Optimización Operativa de Sistemas Agrovoltaicos Dinámicos
%  Equipo: Chacorocalfarobar  |  Organización: Sostenibilidad y Ciencia
% =============================================================================
%
%  INSTRUCCIONES PARA FIGURAS
%  ──────────────────────────
%  Crea la carpeta  sprint3/figures/  y guarda en ella estas capturas de pantalla
%  ejecutando el dashboard con:  streamlit run app.py
%
%  fig01_dashboard_principal.png
%    → Captura completa del Tab 1 ("Dashboard") con el slider a las 12:00,
%      mostrando el SVG solar, el gauge IEC, las cuatro tarjetas de métricas
%      (ePAR S1, ePAR S2, VWC S1, VWC S2) y la caja de justificación en verde.
%
%  fig02_svg_slider.png
%    → Detalle de la columna izquierda del Tab 1: SVG solar con el sol en lo
%      alto, el arco diario, los paneles inclinados y la leyenda de ángulos
%      (azul = actual, verde discontinuo = recomendado).
%
%  fig03_gauge_metricas.png
%    → Detalle de la columna derecha del Tab 1: gauge IEC con las zonas de
%      color y las cuatro tarjetas de ángulo actual, ángulo recomendado,
%      régimen activo y elevación solar.
%
%  fig04_trackers_alertas.png
%    → Sección inferior del Tab 1: grid 2×5 de trackers con sus varianzas y
%      estados (Normal en verde, Alta varianza en rojo), y el panel de alertas
%      activas a la derecha.
%
%  fig05_series_temporales.png
%    → Tab 2 ("Series temporales") con ePAR S1 y VWC S1 seleccionados, mostrando
%      el gráfico de líneas principal con la línea de umbral crítico de VWC (20.0)
%      y los dos gráficos inferiores (tendencia diaria y distribución boxplot).
%
%  fig06_recomendacion.png
%    → Tab 3 ("Recomendación") con la regla activa resaltada en verde y el
%      resto de reglas candidatas listadas por IEC mediana decreciente.
%
%  fig07_alertas.png
%    → Tab 4 ("Alertas") mostrando las alertas activas ordenadas por severidad,
%      la tabla de diagnóstico de trackers con varianzas y el gráfico de
%      evolución reciente de VWC con la línea de umbral crítico.
%
% =============================================================================

\documentclass[11pt,a4paper]{article}

\usepackage[spanish]{babel}
\usepackage[utf8]{inputenc}
\usepackage[T1]{fontenc}
\usepackage{geometry}
\usepackage{longtable}
\usepackage{array}
\usepackage{enumitem}
\usepackage{xcolor}
\usepackage{hyperref}
\usepackage{graphicx}
\usepackage{float}
\usepackage{fancyhdr}
\usepackage{booktabs}
\usepackage{caption}
\usepackage{subcaption}
\usepackage{titlesec}
\usepackage{parskip}

\geometry{margin=2.2cm, top=2.8cm, headheight=15pt}
\setlist[itemize]{label=\textbullet, leftmargin=1.2cm}
\renewcommand{\arraystretch}{1.3}
\sloppy

% ─── Encabezado ───────────────────────────────────────────────────────────────
\pagestyle{fancy}
\fancyhf{}
\fancyhead[L]{\small Sostenibilidad y Ciencia \textperiodcentered{} SAMO}
\fancyhead[R]{\small E05 — Informe de Dashboard v1}
\fancyfoot[C]{\thepage}
\renewcommand{\headrulewidth}{0.4pt}

% ─── Hipervínculos ────────────────────────────────────────────────────────────
\hypersetup{colorlinks=true, linkcolor=black, urlcolor=blue, citecolor=black}

% ─── Captions ─────────────────────────────────────────────────────────────────
\captionsetup{font=small, labelfont=bf, skip=4pt}

% ─── Colores institucionales ──────────────────────────────────────────────────
\definecolor{samogreen}{HTML}{16a34a}
\definecolor{samoblue}{HTML}{2563eb}
\definecolor{samored}{HTML}{dc2626}
\definecolor{samogray}{HTML}{64706d}

% =============================================================================
\begin{document}
\thispagestyle{empty}

% ─── PORTADA ──────────────────────────────────────────────────────────────────
\vspace*{1.2cm}
\begin{center}
  {\large Sostenibilidad y Ciencia}\\[0.6em]
  {\LARGE \textbf{Análisis y Optimización Operativa}}\\
  {\LARGE \textbf{de Sistemas Agrovoltaicos Dinámicos}}\\[2em]
  {\Large \textbf{E05 — Informe de Dashboard v1}}\\[0.5em]
  {\normalsize Sistema de monitorización, recomendación y alertas operativas}\\[2.5em]

  \begin{tabular}{ll}
    \textbf{Proyecto}       & SAMO — Sistemas Agrovoltaicos Dinámicos \\
    \textbf{Equipo}         & Chacorocalfarobar \\
    \textbf{Director PM}    & Daniel Álvarez Sarroca \\
    \textbf{Sprint}         & 3 de 4 \\
    \textbf{Entregable}     & E05 — Dashboard v1 \\
    \textbf{Fecha}          & 4 de mayo de 2026 \\
    \textbf{Versión}        & 1.0 \\
    \textbf{Estado técnico} & Basado en \texttt{main} tras PR~\#2; commit \texttt{2172d58} \\
  \end{tabular}\\[2em]

  \small\textit{Documento técnico interno — prototipo funcional pendiente de validación operativa y agronómica final}
\end{center}

\newpage
\tableofcontents
\newpage

% =============================================================================
\section{Resumen Ejecutivo}

Este informe documenta el trabajo técnico realizado durante el Sprint 3 del proyecto
\textit{Análisis y Optimización Operativa de Sistemas Agrovoltaicos Dinámicos} (SAMO), cuyo
objetivo era construir la primera versión funcional del dashboard operativo. El resultado
principal es un prototipo Streamlit que transforma los outputs del Sprint 2 en una interfaz
de monitorización, recomendación y alerta orientada al operador de la instalación agrovoltaica.

El dashboard opera sobre dos datasets principales: 337 observaciones de modelización a
resolución de 6 horas (base del sistema de recomendación) y 1.462 registros integrados
(base de las series temporales históricas). El IEC observado se mueve entre 0.063 y 0.925,
con media de 0.397, y es el indicador central que el sistema presenta al operador para
evaluar el equilibrio entre producción energética y condiciones microclimáticas del cultivo.

Los resultados más relevantes del sprint son:

\begin{itemize}
  \item Se ha construido un dashboard con cuatro pestañas funcionales: \textbf{Dashboard},
    \textbf{Series temporales}, \textbf{Recomendación} y \textbf{Alertas}.
  \item La pantalla principal integra un slider horario con reproducción automática,
    una visualización SVG del sistema solar-panel, un gauge IEC, tarjetas de
    ePAR/VWC con etiquetas contextuales, un grid de estado de los 10 trackers y un
    resumen de alertas activas.
  \item La arquitectura separa la lógica de negocio (motores de reglas, alertas y cálculo
    solar) de la capa de presentación, con carga cacheada (\texttt{st.cache\_data}) y
    renderizado parcial (\texttt{st.fragment}) para mantener fluidez en la interacción.
  \item Se han corregido dos errores técnicos de alto impacto antes del cierre: la escala
    VWC (umbral 20.0 porcentaje, no 0.20 decimal) y la activación de reglas por condiciones
    reales en lugar de proximidad de IEC.
  \item Los 37 tests unitarios pasan sin errores sobre los módulos críticos del dashboard.
  \item La rama \texttt{iec-sensitivity-dashboard} amplía el sistema con tres prioridades
    de ponderación del IEC (equilibrada, energética y agrícola) y una política de rotación
    por hora para cada escenario.
\end{itemize}

% =============================================================================
\section{Introducción}

\subsection{Contexto del proyecto}

Los sistemas agrovoltaicos dinámicos combinan generación fotovoltaica y cultivo agrícola
en el mismo espacio físico. A diferencia de los sistemas estáticos, los paneles disponen de
un eje de rotación que permite seguir al sol a lo largo del día. Esta rotación afecta
simultáneamente a la captación de irradiancia y al microclima bajo las placas: la sombra
proyectada modula la temperatura del suelo, la humedad volumétrica (VWC) y la radiación
fotosintéticamente activa (ePAR) disponible para el cultivo.

El objetivo del proyecto no es maximizar la producción energética de forma aislada. El
Índice Energía–Cultivo (IEC), definido en el Sprint 2 como combinación ponderada de un
\textit{energy\_score} y un \textit{crop\_score}, cuantifica el equilibrio operativo entre ambas
dimensiones. El dashboard del Sprint 3 hace visible ese equilibrio en tiempo real y
proporciona recomendaciones operativas accionables al operador.

\subsection{Objetivos del Sprint 3}

\begin{itemize}
  \item Diseñar e implementar el Dashboard v1 como entregable E05, sobre los outputs del Sprint 2.
  \item Mostrar el estado actual o simulado de los sensores clave mediante un slider horario.
  \item Visualizar series temporales filtrables por variable, tipo y rango de fechas.
  \item Integrar la recomendación de la regla de rotación aplicable, con evaluación de
    condiciones reales del registro activo.
  \item Detectar y notificar alertas operativas: anomalías angulares de trackers y tendencias
    críticas de humedad del suelo.
  \item Garantizar la calidad del código con una suite de tests unitarios antes del cierre.
  \item Documentar decisiones técnicas, limitaciones y próximos pasos para el Sprint 4.
\end{itemize}

\subsection{Alcance del informe}

El informe cubre el entregable técnico E05 y la implementación actual del directorio
\texttt{sprint3/}. No sustituye el documento de gestión del sprint (\texttt{Gestion\_Proyecto\_Sprint3.tex}),
pero describe con detalle los componentes del dashboard, las decisiones de diseño, los
resultados observados y los riesgos técnicos relevantes para la entrega.

% =============================================================================
\section{Arquitectura del Dashboard}

\subsection{Diseño general}

La solución está implementada en Python con Streamlit y separa estrictamente la lógica
de negocio de la capa de presentación. \texttt{app.py} actúa como punto de entrada: configura
la página, carga los cuatro datasets con caché y delega la visualización a los módulos de
tab. Los motores de reglas, alertas, cálculo solar y generación SVG se mantienen en módulos
independientes para facilitar las pruebas unitarias (Figura~\ref{fig:arquitectura}).

\begin{longtable}{p{0.32\textwidth} p{0.6\textwidth}}
\toprule
\textbf{Componente} & \textbf{Responsabilidad} \\
\midrule
\texttt{app.py} & Configuración de la página, carga cacheada de datos y estructura de las cuatro pestañas \\
\texttt{data\_loader.py} & Carga de los cuatro CSVs del Sprint 2 con \texttt{st.cache\_data}; ordenación temporal \\
\texttt{tabs/tab\_estado.py} & Tab 1: slider horario, SVG solar, gauge IEC, tarjetas de métricas, grid de trackers y alertas \\
\texttt{tabs/tab\_series.py} & Tab 2: series temporales con Plotly, filtros, tendencia diaria y distribución boxplot \\
\texttt{tabs/tab\_recomendacion.py} & Tab 3: política de rotación y regla activa con evaluación de condiciones reales \\
\texttt{tabs/tab\_alertas.py} & Tab 4: alertas ordenadas por severidad, tabla de trackers y evolución VWC \\
\texttt{rule\_engine.py} & Motor de evaluación de las 6 reglas candidatas del Sprint 2 \\
\texttt{alert\_logic.py} & Detección de anomalías angulares de trackers y tendencia VWC \\
\texttt{solar\_logic.py} & Estimación de elevación solar y ángulo óptimo recomendado por hora \\
\texttt{svg\_generator.py} & Generación del SVG animado con sol, arco diario, paneles y ángulos \\
\texttt{styles.py} & Inyección de CSS global, paleta de colores y componentes HTML reutilizables \\
\bottomrule
\end{longtable}

\subsection{Flujo de datos}

\begin{enumerate}
  \item Los cuatro CSVs del Sprint 2 se leen desde \texttt{sprint2/outputs\_sprint2/} al arrancar
    la aplicación. \texttt{st.cache\_data} evita relecturas en cada interacción posterior.
  \item El usuario selecciona una hora mediante el slider (rango 06:00–21:00). El sistema
    busca el registro con mayor IEC para esa hora; si no existe, aplica un fallback al
    registro más próximo disponible.
  \item El módulo \texttt{solar\_logic.py} calcula la elevación solar teórica y el ángulo óptimo
    de referencia para la hora seleccionada.
  \item \texttt{rule\_engine.py} evalúa las condiciones reales del registro activo (Albedo, hora,
    régimen, VWC, GPOA) y determina qué regla candidata, si alguna, está activa.
  \item \texttt{alert\_logic.py} genera la lista de alertas activas comparando la varianza angular
    de cada tracker con el umbral de 450\,deg\textsuperscript{2} y la tendencia reciente de VWC
    con el umbral crítico de 20.0\,pp.
  \item La UI renderiza todos los componentes de la vista activa y el header del dashboard
    muestra el número total de alertas en tiempo real.
\end{enumerate}

\subsection{Decisiones de rendimiento y UX}

\begin{itemize}
  \item \textbf{\texttt{st.cache\_data}}: evita relecturas de CSV en cada interacción de usuario.
    Los datasets son estáticos; la caché se invalida únicamente al reiniciar el servidor.
  \item \textbf{\texttt{st.fragment(run\_every=``3s'')}}: el slider horario y la sección interactiva
    principal del Tab 1 se envuelven en un fragmento que se rerenderiza cada 3 segundos
    cuando está activo el modo de reproducción automática. La sección de trackers y alertas
    permanece fuera del fragmento para evitar parpadeos.
  \item \textbf{SVG renderizado como imagen base64}: el SVG solar se codifica en base64 y se
    embebe como etiqueta \texttt{<img>}. Streamlit sanitiza el HTML de \texttt{st.markdown} eliminando
    defs, gradientes y filters; el método base64 conserva todos los estilos SVG.
  \item \textbf{Paleta visual}: verde \texttt{\#16a34a} como color operativo principal, azul
    \texttt{\#2563eb} para valores de referencia y rojo \texttt{\#dc2626} para alertas. La tipografía
    es Apple/System para coherencia visual con el sistema operativo del cliente.
  \item \textbf{Reproducción automática}: el slider avanza automáticamente un paso por hora
    cada 3 segundos mediante un toggle de control. El modo puede desactivarse manualmente
    o via parámetro URL (\texttt{?autoplay=0}).
\end{itemize}

% =============================================================================
\section{Datos de Entrada}

\begin{longtable}{p{0.40\textwidth} p{0.40\textwidth} p{0.12\textwidth}}
\toprule
\textbf{Archivo} & \textbf{Uso en Sprint 3} & \textbf{Registros} \\
\midrule
\texttt{dataset\_modelizacion\_6h.csv} & IEC, régimen, tracking, ePAR/VWC medios, elevación solar — base del Tab 1 y Tab 3 & 337 \\
\texttt{dataset\_integrado\_6h.csv} & Series temporales históricas de todas las variables — base del Tab 2 & 1.462 \\
\texttt{candidate\_rotation\_rules.csv} & Las 6 reglas candidatas de política de rotación — Tab 1 y Tab 3 & 6 \\
\texttt{tracker\_variance\_diagnostic.csv} & Varianza angular por tracker — Tab 1 (grid) y Tab 4 (diagnóstico) & 10 \\
\bottomrule
\end{longtable}

\subsection{Estadísticos de las variables principales}

\begin{longtable}{p{0.25\textwidth} p{0.13\textwidth} p{0.13\textwidth} p{0.13\textwidth} p{0.28\textwidth}}
\toprule
\textbf{Variable} & \textbf{Mínimo} & \textbf{Media} & \textbf{Máximo} & \textbf{Interpretación} \\
\midrule
IEC                 & 0.063  & 0.397   & 0.925    & Equilibrio energía-cultivo (0–1) \\
VWC\_S1\_mean       & 10.405 & 22.786  & 39.700   & Humedad suelo S1 (\% volumétrico) \\
ePAR\_S1\_mean      & $-0.016$ & 444.754 & 2.319.500 & Radiación PAR (\textmu mol/m\textsuperscript{2}/s) \\
track\_mean         & $-32.3$ & 0.034  & 33.100   & Ángulo medio de seguimiento (°) \\
solar\_elevation\_deg & $-18.2$ & 25.4  & 71.3    & Elevación solar teórica (°) \\
\bottomrule
\end{longtable}

\subsection{Distribución por régimen de operación}

\begin{longtable}{p{0.35\textwidth} p{0.2\textwidth} p{0.2\textwidth} p{0.18\textwidth}}
\toprule
\textbf{Régimen} & \textbf{Observaciones} & \textbf{IEC media} & \textbf{IEC mediana} \\
\midrule
TRACKING\_PM  & 84  & 0.742 & 0.741 \\
TRACKING\_AM  & 84  & 0.440 & 0.429 \\
HORIZONTAL    & 169 & 0.204 & 0.184 \\
\bottomrule
\end{longtable}

HORIZONTAL representa el 50.1\% de las observaciones y tiene el IEC más bajo
(mediana 0.184). La diferencia entre TRACKING\_PM y HORIZONTAL es de 0.557 puntos de
IEC mediana, umbral que el dashboard usa para las alertas de régimen.

% =============================================================================
\section{Funcionalidades Implementadas}

\subsection{Tab 1 — Dashboard principal}

La pantalla principal concentra el estado operativo completo del sistema en una única vista
(Figura~\ref{fig:dashboard_principal}). El slider horario (rango 06:00–21:00) permite al
operador explorar cualquier franja del día y ver cómo cambian el régimen, el ángulo y el
IEC. La reproducción automática avanza el slider cada 3 segundos, simulando una visión
de evolución diaria del sistema.

\begin{figure}[H]
  \centering
  % Guarda la captura en sprint3/figures/fig01_dashboard_principal.png
  \includegraphics[width=\textwidth]{figures/fig01_dashboard_principal.png}
  \caption{Tab 1 — Dashboard principal. Slider a las 12:00 (TRACKING\_PM). De arriba a abajo:
    caja de justificación en lenguaje natural, SVG solar (izquierda) con gauge IEC y tarjetas
    de métricas (derecha), y las 6 reglas candidatas con la regla activa resaltada en verde.}
  \label{fig:dashboard_principal}
\end{figure}

Componentes del Tab 1:

\begin{itemize}
  \item \textbf{Caja de justificación en lenguaje natural}: para cada hora y régimen, el sistema
    genera un texto contextual que explica por qué los paneles están en esa posición, la
    elevación solar en ese instante y el nivel del IEC. El texto varía entre TRACKING\_PM
    (fondo verde, explicación de máximo rendimiento), TRACKING\_AM (fondo azul) y
    HORIZONTAL (fondo neutro, razón de reposo o protección de viento).
  \item \textbf{SVG solar} (Figura~\ref{fig:svg_slider}): visualización vectorial con el sol
    posicionado en el arco diario, los paneles inclinados al ángulo actual (línea azul sólida)
    y la referencia del ángulo óptimo (línea verde discontinua). Incluye etiquetas de ángulo
    y un indicador de estado «En rango óptimo» / «Fuera del rango recomendado».
  \item \textbf{Gauge IEC y tarjetas de métricas} (Figura~\ref{fig:gauge_metricas}): el gauge
    Plotly muestra el IEC con tres zonas coloreadas: roja (crítico $<$0.35), amarilla
    (media 0.35–0.60) y verde (óptimo $>$0.60). Las cuatro tarjetas muestran ángulo actual,
    ángulo recomendado, régimen activo y elevación solar.
  \item \textbf{Tarjetas ePAR y VWC}: cuatro cards con etiquetas contextuales. ePAR evalúa si
    la irradiancia PAR supera el umbral de compensación lumínica (200\,\textmu mol/m\textsuperscript{2}/s).
    VWC evalúa si la humedad supera el umbral de riesgo hídrico (20.0\,pp) o está bien
    hidratada ($\geq$30.0\,pp).
  \item \textbf{Grid de 10 trackers} (Figura~\ref{fig:trackers_alertas}): cada tracker muestra
    su varianza angular en deg\textsuperscript{2} y su estado (Normal en verde si $\leq$450\,deg\textsuperscript{2},
    Alta varianza en rojo si $>$450\,deg\textsuperscript{2}).
  \item \textbf{Panel de alertas activas}: a la derecha del grid, lista de alertas por tipo
    (tracker anómalo, tendencia VWC crítica) con descripción y nivel de severidad.
\end{itemize}

\begin{figure}[H]
  \centering
  \begin{subfigure}[t]{0.54\textwidth}
    % Guarda en sprint3/figures/fig02_svg_slider.png
    \includegraphics[width=\textwidth]{figures/fig02_svg_slider.png}
    \caption{SVG solar con el ángulo actual (azul) y el óptimo recomendado (verde discontinuo).}
    \label{fig:svg_slider}
  \end{subfigure}
  \hfill
  \begin{subfigure}[t]{0.43\textwidth}
    % Guarda en sprint3/figures/fig03_gauge_metricas.png
    \includegraphics[width=\textwidth]{figures/fig03_gauge_metricas.png}
    \caption{Gauge IEC y tarjetas de ángulo, régimen y elevación solar.}
    \label{fig:gauge_metricas}
  \end{subfigure}
  \caption{Detalle de la columna visual del Tab 1 para la franja de las 12:00 (TRACKING\_PM, IEC~0.925).}
\end{figure}

\begin{figure}[H]
  \centering
  % Guarda en sprint3/figures/fig04_trackers_alertas.png
  \includegraphics[width=\textwidth]{figures/fig04_trackers_alertas.png}
  \caption{Sección inferior del Tab 1: grid de estado de los 10 trackers (M02, M05, M06, M07, M10 en rojo)
    y panel de alertas activas con severidad.}
  \label{fig:trackers_alertas}
\end{figure}

\subsection{Tab 2 — Series temporales}

El Tab 2 permite al operador explorar la evolución histórica de las variables del sistema
(Figura~\ref{fig:series_temporales}). El selector de variables ofrece 13 columnas agrupadas
en cuatro categorías: ePAR (S1/S2, dos sensores por sección), VWC (S1/S2, dos sensores por
sección), irradiancia GPOA (S1/S2) y ángulo de tracking (M01, M03, M05).

El tab genera tres visualizaciones complementarias:

\begin{enumerate}
  \item \textbf{Gráfico de líneas principal}: evolución temporal de las variables seleccionadas
    sobre el rango de fechas filtrado. Las líneas de umbral crítico se superponen en rojo
    discontinuo (ePAR: 200\,\textmu mol/m\textsuperscript{2}/s) y en naranja discontinuo (VWC: 20.0\,pp).
  \item \textbf{Tendencia diaria suavizada}: gráfico de área con la media diaria de cada variable,
    que permite identificar patrones de semanas o meses sin el ruido de la resolución de 6 horas.
  \item \textbf{Distribución de valores} (boxplot): distribución estadística de cada variable
    en el rango seleccionado, útil para comparar el comportamiento de S1 frente a S2 o
    identificar asimetrías y outliers.
\end{enumerate}

Cuatro tarjetas resumen muestran la ventana temporal activa, el número de valores disponibles,
el pico máximo de ePAR y el punto hídrico más bajo de VWC en el período filtrado.

\begin{figure}[H]
  \centering
  % Guarda en sprint3/figures/fig05_series_temporales.png
  \includegraphics[width=\textwidth]{figures/fig05_series_temporales.png}
  \caption{Tab 2 — Series temporales con ePAR S1 y VWC S1 seleccionados. La línea naranja
    discontinua marca el umbral crítico de VWC (20.0\,pp). Los dos gráficos inferiores muestran
    la tendencia diaria (área) y la distribución estadística (boxplot).}
  \label{fig:series_temporales}
\end{figure}

\subsection{Tab 3 — Recomendación}

El Tab 3 presenta la política de rotación activa y la regla candidata aplicable al registro
activo (Figura~\ref{fig:recomendacion}). A diferencia de la versión inicial del Sprint 3 (que
seleccionaba la regla por proximidad de IEC), el motor actual evalúa las condiciones reales
del registro más reciente: Albedo\_S1, Albedo\_S2, elevación solar, hora del día, régimen
activo, VWC\_S1 y GPOA\_S2. Si ninguna condición se satisface, no se marca ninguna regla
como activa, evitando recomendaciones espurias.

Las 6 reglas candidatas se presentan ordenadas por IEC mediana esperado (de mayor a menor),
con el tipo de regla en un badge de color, el soporte estadístico y el texto de la condición.
La regla activa se resalta con un borde verde y una etiqueta «ACTIVA».

\begin{figure}[H]
  \centering
  % Guarda en sprint3/figures/fig06_recomendacion.png
  \includegraphics[width=\textwidth]{figures/fig06_recomendacion.png}
  \caption{Tab 3 — Recomendación. La regla de prioridad alta (Albedo\_S1 $>$ 55.7 y elevación
    $>$ 68°, IEC mediana 0.862) aparece marcada como ACTIVA con borde verde. Las cinco
    reglas restantes se listan con su soporte estadístico.}
  \label{fig:recomendacion}
\end{figure}

\subsection{Tab 4 — Alertas}

El Tab 4 centraliza la información de anomalías del sistema (Figura~\ref{fig:alertas}). Incluye
tres componentes:

\begin{itemize}
  \item \textbf{Lista de alertas activas}: ordenada por severidad (CRÍTICO en rojo, AVISO en
    naranja). Cada alerta incluye título, descripción detallada y nivel.
  \item \textbf{Tabla de diagnóstico de trackers}: todos los trackers con su varianza angular
    en deg\textsuperscript{2}, ordenados de mayor a menor varianza. Los trackers que superan el
    umbral de 450\,deg\textsuperscript{2} se destacan como de Alta varianza.
  \item \textbf{Evolución reciente de VWC}: gráfico de evolución temporal de VWC en los
    últimos registros disponibles, con la línea de umbral crítico (20.0\,pp) superpuesta.
    Las alertas de VWC se expresan como tendencia en puntos porcentuales por hora.
\end{itemize}

\begin{figure}[H]
  \centering
  % Guarda en sprint3/figures/fig07_alertas.png
  \includegraphics[width=\textwidth]{figures/fig07_alertas.png}
  \caption{Tab 4 — Alertas. Panel superior con las alertas activas por severidad. Tabla inferior
    izquierda con el diagnóstico de varianza de los 10 trackers. Gráfico inferior derecho con la
    evolución reciente de VWC y la línea de umbral crítico.}
  \label{fig:alertas}
\end{figure}

% =============================================================================
\section{Reglas, Alertas y Control Operativo}

\subsection{Reglas candidatas de rotación}

Las 6 reglas candidatas proceden del análisis del Sprint 2 (árbol de decisión + análisis
observacional por régimen) y están implementadas en \texttt{rule\_engine.py} con evaluación
de condiciones explícitas.

\begin{longtable}{p{0.16\textwidth} p{0.45\textwidth} p{0.14\textwidth} p{0.1\textwidth} p{0.07\textwidth}}
\toprule
\textbf{Tipo} & \textbf{Condición / recomendación} & \textbf{Condición evaluada} & \textbf{IEC mediana} & \textbf{n} \\
\midrule
prioridad\_alta & Si Albedo\_S1 $>$ 55.7 y la elevación solar supera 68°, priorizar tracking activo de tarde y evitar modos conservadores & Albedo\_S1, solar\_elev & 0.862 & 33 \\
operativa & Si la franja es 12:00 y el sistema opera en TRACKING\_PM, mantener un ángulo cercano a 31.8° & hora, régimen & 0.769 & 75 \\
operativa & Si la franja es 06:00 y el sistema opera en TRACKING\_AM, mantener un ángulo cercano a $-$32.2° & hora, régimen & 0.602 & 10 \\
transicion & Si Albedo\_S1 $\leq$ 55.7 pero Albedo\_S2 $>$ 18.0 y VWC\_S1 $>$ 20.0, usar tracking suave antes que HORIZONTAL permanente & Albedo, VWC\_S1 & 0.435 & 51 \\
evitacion & Evitar HORIZONTAL en franjas productivas cuando existan condiciones para tracking activo & (umbral régimen) & 0.184 & 169 \\
cautela\_hidrica & Si Albedo\_S1 $\leq$ 55.7, Albedo\_S2 $\leq$ 18.0, VWC\_S1 $\leq$ 22.2 y GPOA\_S2 $\leq$ 35.5, evitar forzar producción: IEC esperado muy bajo & Albedo, VWC, GPOA & 0.093 & 43 \\
\bottomrule
\end{longtable}

La corrección implementada en el PR~\#2 garantiza que la activación de reglas evalúa las
columnas reales del registro activo. Si el registro no contiene alguna variable necesaria,
la regla correspondiente no se activa. Esto elimina los falsos positivos que se producían
en la versión inicial, donde la regla con IEC más próximo al IEC actual se marcaba como activa
con independencia de sus condiciones.

\subsection{Diagnóstico de trackers}

El umbral de anomalía se fija en 450\,deg\textsuperscript{2}. Una varianza superior indica
que el tracker no sigue con precisión las consignas angulares, lo que puede deberse a un
fallo mecánico, una posición de stow permanente o un problema de comunicación.

\begin{longtable}{p{0.2\textwidth} p{0.26\textwidth} p{0.2\textwidth} p{0.28\textwidth}}
\toprule
\textbf{Tracker} & \textbf{Varianza (deg\textsuperscript{2})} & \textbf{Estado} & \textbf{Nota} \\
\midrule
M01 & 445.2 & Normal  & Por debajo del umbral \\
M02 & 508.5 & \textcolor{samored}{Alta varianza} & Identificado como anómalo en Sprint 1 \\
M03 & 441.7 & Normal  & \\
M04 & 438.1 & Normal  & \\
M05 & 534.9 & \textcolor{samored}{Alta varianza} & Varianza más alta del conjunto \\
M06 & 513.0 & \textcolor{samored}{Alta varianza} & Identificado como anómalo en Sprint 1 \\
M07 & 535.0 & \textcolor{samored}{Alta varianza} & Varianza más alta del conjunto \\
M08 & 443.5 & Normal  & \\
M09 & 446.8 & Normal  & \\
M10 & 509.0 & \textcolor{samored}{Alta varianza} & Identificado como anómalo en Sprint 1 \\
\bottomrule
\end{longtable}

Cinco de los diez trackers superan el umbral: M02, M05, M06, M07 y M10. M02, M06 y M10
fueron identificados con comportamiento anómalo ya en el Sprint 1 (posición prácticamente
fija en torno a 50.6°). M05 y M07 son trackers activos que no estaban en la lista original
del Sprint 1, por lo que sus varianzas altas merecen verificación técnica adicional.

\subsection{Lógica de alertas}

El módulo \texttt{alert\_logic.py} genera alertas de dos tipos:

\begin{itemize}
  \item \textbf{Alertas de tracker}: se genera una alerta por cada tracker con varianza angular
    superior a 450\,deg\textsuperscript{2}. La severidad es CRÍTICO para varianzas $>$520\,deg\textsuperscript{2}
    y AVISO para el rango 450–520\,deg\textsuperscript{2}.
  \item \textbf{Alertas de VWC}: se evalúa la tendencia reciente de VWC en los últimos registros
    del dataset de modelización. Una caída sostenida que sitúe la proyección por debajo de
    20.0\,pp genera una alerta de riesgo hídrico con la tendencia expresada en puntos
    porcentuales por hora.
\end{itemize}

% =============================================================================
\section{Mejora Técnica: Sensibilidad del IEC por Prioridades}

La rama \texttt{iec-sensitivity-dashboard} (no fusionada en \texttt{main} en esta versión) trata
el IEC como índice configurable en lugar de verdad fija. Permite comparar tres escenarios
de ponderación que responden a diferentes estrategias operativas del cliente.

\begin{longtable}{p{0.22\textwidth} p{0.18\textwidth} p{0.18\textwidth} p{0.36\textwidth}}
\toprule
\textbf{Escenario} & \textbf{Peso energía} & \textbf{Peso cultivo} & \textbf{Objetivo operativo} \\
\midrule
Equilibrado       & 0.5 & 0.5 & Compromiso simétrico entre producción y cultivo \\
Priorizar energía & 0.7 & 0.3 & Maximizar la generación fotovoltaica \\
Priorizar cultivo & 0.3 & 0.7 & Favorecer la protección microclimática y hídrica \\
\bottomrule
\end{longtable}

\subsection{Resultados por régimen y escenario}

\begin{longtable}{p{0.18\textwidth} p{0.18\textwidth} p{0.08\textwidth} p{0.15\textwidth} p{0.15\textwidth} p{0.12\textwidth} p{0.12\textwidth}}
\toprule
\textbf{Escenario} & \textbf{Régimen} & \textbf{n} & \textbf{IEC orig.\ med.} & \textbf{IEC pond.\ med.} & \textbf{Energy med.} & \textbf{Crop med.} \\
\midrule
Equilibrado       & HORIZONTAL    & 169 & 0.184 & 0.184 & 0.000 & 0.264 \\
Equilibrado       & TRACKING\_AM  & 84  & 0.429 & 0.429 & 0.442 & 0.485 \\
Equilibrado       & TRACKING\_PM  & 84  & 0.741 & 0.741 & 0.908 & 0.626 \\
\midrule
Priorizar energía & HORIZONTAL    & 169 & 0.184 & 0.142 & 0.000 & 0.264 \\
Priorizar energía & TRACKING\_AM  & 84  & 0.429 & 0.419 & 0.442 & 0.485 \\
Priorizar energía & TRACKING\_PM  & 84  & 0.741 & 0.811 & 0.908 & 0.626 \\
\midrule
Priorizar cultivo & HORIZONTAL    & 169 & 0.184 & 0.208 & 0.000 & 0.264 \\
Priorizar cultivo & TRACKING\_AM  & 84  & 0.429 & 0.444 & 0.442 & 0.485 \\
Priorizar cultivo & TRACKING\_PM  & 84  & 0.741 & 0.695 & 0.908 & 0.626 \\
\bottomrule
\end{longtable}

\subsection{Política recomendada por hora}

\begin{longtable}{p{0.20\textwidth} p{0.08\textwidth} p{0.22\textwidth} p{0.12\textwidth} p{0.18\textwidth} p{0.10\textwidth}}
\toprule
\textbf{Escenario} & \textbf{Hora} & \textbf{Régimen rec.} & \textbf{Ángulo (°)} & \textbf{IEC pond.\ med.} & \textbf{Soporte} \\
\midrule
Equilibrado       & 00 & HORIZONTAL   & $-2.7$  & 0.300 & 84 \\
Equilibrado       & 06 & TRACKING\_AM & $-32.3$ & 0.613 & 84 \\
Equilibrado       & 12 & TRACKING\_PM & $+32.6$ & 0.925 & 84 \\
Equilibrado       & 18 & HORIZONTAL   & $+2.3$  & 0.430 & 85 \\
\midrule
Priorizar energía & 00 & HORIZONTAL   & $-2.7$  & 0.180 & 84 \\
Priorizar energía & 06 & TRACKING\_AM & $-31.5$ & 0.716 & 84 \\
Priorizar energía & 12 & TRACKING\_PM & $+32.6$ & 0.955 & 84 \\
Priorizar energía & 18 & HORIZONTAL   & $+2.3$  & 0.362 & 85 \\
\midrule
Priorizar cultivo & 00 & HORIZONTAL   & $-2.7$  & 0.421 & 84 \\
Priorizar cultivo & 06 & TRACKING\_AM & $-32.3$ & 0.667 & 84 \\
Priorizar cultivo & 12 & TRACKING\_PM & $+32.6$ & 0.895 & 84 \\
Priorizar cultivo & 18 & HORIZONTAL   & $+2.3$  & 0.498 & 85 \\
\bottomrule
\end{longtable}

El análisis confirma que TRACKING\_PM es el régimen más sólido en todos los escenarios
para la franja de las 12:00. Al priorizar energía, su IEC ponderado sube de 0.741 a 0.811
(+9.4\%); al priorizar cultivo baja a 0.695 ($-$6.2\%), pero sigue siendo competitivo frente
a las alternativas. HORIZONTAL mejora relativamente cuando se prioriza cultivo (de 0.184 a
0.208 en mediana), pero no supera a los modos de seguimiento activo en ninguna franja productiva.
La franja de las 18:00 es el caso donde la elección de ponderación tiene mayor impacto
relativo sobre la recomendación: el IEC de HORIZONTAL varía entre 0.362 (energía) y 0.498
(cultivo), un rango de 0.136 puntos que puede ser operativamente significativo.

% =============================================================================
\section{Validación Técnica}

\subsection{Suite de tests unitarios}

El Sprint 3 incorporó por primera vez una suite de tests unitarios con cobertura de todos
los módulos de lógica de negocio. Los tests se organizan en seis archivos en el directorio
\texttt{sprint3/tests/}:

\begin{longtable}{p{0.38\textwidth} p{0.15\textwidth} p{0.4\textwidth}}
\toprule
\textbf{Archivo de test} & \textbf{Tests} & \textbf{Módulo cubierto} \\
\midrule
\texttt{test\_alert\_logic.py}    & 7  & \texttt{alert\_logic.py}: detección de trackers anómalos y tendencia VWC \\
\texttt{test\_data\_loader.py}    & 5  & \texttt{data\_loader.py}: carga, ordenación y caché de CSVs \\
\texttt{test\_rule\_engine.py}    & 9  & \texttt{rule\_engine.py}: activación y no activación de cada regla \\
\texttt{test\_solar\_logic.py}    & 6  & \texttt{solar\_logic.py}: elevación solar y ángulo recomendado \\
\texttt{test\_svg\_generator.py}  & 5  & \texttt{svg\_generator.py}: generación de SVG válido y parámetros \\
\texttt{test\_vwc\_thresholds.py} & 5  & Umbrales VWC (20.0 / 30.0) en lógica de etiquetado \\
\midrule
\textbf{Total} & \textbf{37} & \\
\bottomrule
\end{longtable}

Comando de ejecución: \texttt{py -3.11 -m pytest tests -q -p no:cacheprovider}

Resultado: \textbf{37 passed} — sin errores ni warnings en el conjunto de test actual.

Los tests de \texttt{test\_rule\_engine.py} y \texttt{test\_vwc\_thresholds.py} fueron los que
detectaron los dos bugs corregidos en el PR~\#2:

\begin{itemize}
  \item \textbf{Bug de escala VWC}: el test \texttt{test\_vwc\_thresholds.py} verificó que el
    umbral de alerta es 20.0 (porcentaje) y no 0.20 (decimal), lo que habría generado alertas
    de VWC incorrectas en la práctica totalidad de los registros.
  \item \textbf{Bug de activación de regla}: el test \texttt{test\_rule\_engine.py} comprobó que
    ninguna regla se activa cuando sus condiciones no se cumplen, detectando que la versión
    previa al PR marcaba como activa la regla con IEC más próximo al valor actual con
    independencia de las condiciones declaradas.
\end{itemize}

\subsection{Integridad del repositorio}

\begin{longtable}{p{0.48\textwidth} p{0.46\textwidth}}
\toprule
\textbf{Verificación} & \textbf{Resultado} \\
\midrule
37 tests unitarios & \textcolor{samogreen}{\textbf{37 passed}} \\
PR \#2 fusionado en \texttt{main} & Commit \texttt{2172d58} \\
Working tree limpio al cierre & Sin cambios no versionados \\
Corrección escala VWC & Umbrales 20.0 / 30.0 en UI, series y alertas \\
Corrección regla activa & Evaluación de condiciones reales del registro \\
\bottomrule
\end{longtable}

% =============================================================================
\section{Limitaciones y Próximos Pasos}

\subsection{Limitaciones actuales}

\begin{itemize}
  \item \textbf{Sistema de soporte a la decisión, no de control}: el dashboard proporciona
    recomendaciones operativas, no consignas automáticas. Cualquier acción sobre los trackers
    requiere intervención humana validada.
  \item \textbf{Resolución de 6 horas}: el dataset de modelización solo contiene franjas a las
    0, 6, 12 y 18 horas. Las horas intermedias del slider usan el registro con el IEC más
    alto de la hora disponible más cercana, lo que introduce una aproximación en la representación
    del estado del sistema.
  \item \textbf{Reglas candidatas sin validación agronómica}: las 6 reglas proceden del Sprint 2
    y reflejan correlaciones estadísticas, no causalidad probada. Se etiquetan explícitamente
    como provisionales en el dashboard hasta recibir validación del especialista agrícola.
  \item \textbf{Umbrales IEC/VWC/ePAR operativos}: los valores de referencia del IEC (pesos
    0.5/0.5) y los umbrales de VWC (20.0 y 30.0\,pp) son operativos. Deben contrastarse
    con el especialista agrícola del proyecto antes de la entrega final.
  \item \textbf{Rama de sensibilidad IEC no fusionada}: el análisis de ponderaciones es una
    extensión técnica relevante, pero no forma parte de \texttt{main} en esta versión.
  \item \textbf{Datos históricos únicamente}: el dashboard no tiene feed en tiempo real. Opera
    sobre los datos del período junio–octubre 2025.
\end{itemize}

\subsection{Recomendaciones para el Sprint 4}

\begin{itemize}
  \item Decidir si la rama \texttt{iec-sensitivity-dashboard} se integra en \texttt{main} para la
    demo final. Si el cliente valora la comparativa de escenarios, la integración añade valor
    diferencial sin coste técnico adicional relevante.
  \item Capturar evidencias visuales del dashboard funcional (pantallazos de los cuatro tabs)
    para incorporarlas a la documentación técnica final E07 y a la presentación de cierre E06.
  \item Limpiar los artefactos \texttt{\_\_pycache\_\_} versionados en el repositorio y actualizar
    el \texttt{.gitignore} para evitar que vuelvan a versionarse.
  \item Revisar con el especialista agrícola las 6 reglas y los umbrales del IEC antes de la
    presentación final, y documentar el resultado como artefacto del proyecto.
  \item Preparar una guía de ejecución breve (2–3 pasos) para que el cliente pueda lanzar el
    dashboard de forma autónoma en su entorno.
  \item Evaluar si el alcance del Sprint 4 puede incluir un módulo de ingesta en tiempo real
    o si el alcance del proyecto queda explícitamente restringido al análisis histórico.
\end{itemize}

% =============================================================================
\section{Conclusión}

El Sprint 3 ha cumplido su objetivo principal: transformar los outputs del análisis y la
modelización en una interfaz operativa funcional. El dashboard v1 integra en una sola
aplicación la monitorización del IEC, la visualización histórica de sensores, la
recomendación de régimen de rotación y la detección de alertas, siguiendo en todo momento
el principio rector del proyecto: el equilibrio entre producción energética y condiciones
favorables para el cultivo.

Los dos bugs corregidos antes del cierre (escala VWC y activación de reglas) ilustran la
importancia de la suite de tests unitarios como mecanismo de calidad. Sin esa cobertura,
ambos errores habrían llegado a la entrega con impacto directo sobre la fiabilidad de las
alertas y las recomendaciones.

El sistema está listo para ser presentado al cliente en la Sprint Review del Sprint 4 como
prototipo funcional, con las limitaciones explícitamente documentadas y una ruta clara
hacia la versión de entrega final.

\vspace{2em}
\begin{center}
  \textit{Equipo Chacorocalfarobar — Sostenibilidad y Ciencia}
\end{center}

\end{document}
